%%%
% src::
%     urls = https://tex.stackexchange.com/a/654470,
%            https://tex.stackexchange.com/a/761377/6880
%%%

\documentclass[border = 3pt]{standalone}

%%%
% Le \pack ''nicematrix'' est le couteau suisse des \mats.
%%%
\usepackage{nicematrix}

%%%
% Le \pack ''tikz'' sert à décorer le \tab d'adjacence.
%%%
\usepackage{tikz}

\begin{document}

%%%
% La commande \renewcommand{\arraystretch}{1.2} sert à augmenter l'espacement vertical (la hauteur) des lignes dans les tableaux et les matrices en LaTeX.
%%%
\renewcommand{\arraystretch}{1.2}

%%%
% XX
%
% Nous souhaitons tracer \ttes les lignes verticales et horizontales.
%
% On exclut le coin nord-ouest du tableau.
%%%
\begin{NiceTabular}{>{\bfseries}*{6}{c}}[
  hvlines,
  corners = NW
]
\CodeBefore [create-cell-nodes]
  \tikz \draw[red!60, ->, shorten <= 1pt, shorten >= 1pt]
    (4-1.east)
      -| node[pos = 0.5, name = corner] {}
    (1-2.south);
  \tikz \fill[red!10] (corner) circle (5pt);
\Body
  \RowStyle{\bfseries}
    & A & B & C & D & E  \\
  A & 0 & 0 & 1 & 1 & 0  \\
  B & 0 & 0 & 0 & 1 & 1  \\
  C & 1 & 0 & 0 & 0 & 1  \\
  D & 1 & 1 & 0 & 0 & 0  \\
  E & 0 & 1 & 1 & 0 & 0
\end{NiceTabular}

\end{document}
